\subsection{Native points with an arbitrary prescribed coefficient seminorm}
\label{subsec:relative-native-points}

The coefficient ring may be any commutative ring. Its prescribed real
ultrametric seminorm is fixed on every coefficient, including zero.
Only the affine native base point space is known to be Hausdorff; no
Hausdorff hypothesis or conclusion for an arbitrary scheme is used.

\begin{theorem}[Exact closed native coefficient locus]
\label{scheme:relative-valued-points}
\lean{EnvelopeProperty.Berkovich.SchemeValuedPoint.basePoint,
EnvelopeProperty.Berkovich.SchemeValuedPoint.spectrumHomeomorph_val_apply,
EnvelopeProperty.Berkovich.SchemeValuedPoint.basePoint_eval,
EnvelopeProperty.Berkovich.SchemeValuedPoint.relativeLocus,
EnvelopeProperty.Berkovich.SchemeValuedPoint.mem_relativeLocus,
EnvelopeProperty.Berkovich.SchemeValuedPoint.isClosed_relativeLocus,
EnvelopeProperty.Berkovich.RelativeSchemeValuedPoint,
EnvelopeProperty.Berkovich.RelativeSchemeValuedPoint.isClosedEmbedding_subtypeVal}
\leanok
\uses{scheme:affine-valued-point-homeomorphism,scheme:valued-point-pullback,
scheme:valued-point-topology,spectrum:relative}
Let $R$ be a commutative ring with prescribed real seminorm $v$ and let
$f:X\to\operatorname{Spec}R$ be an actual scheme morphism. Let the native
base point be the inverse image of the genuine point $v$ under the native
affine homeomorphism. Native points of $X$ mapping to this exact whole base
point form a closed subspace.
\end{theorem}
\begin{proof}
The proved affine homeomorphism identifies the native point space of
$\operatorname{Spec}R$ with the Hausdorff seminorm spectrum of $R$, so a
singleton native base point is closed. The actual native map induced by
$f$ is continuous, and the exact coefficient locus is its inverse image
of that singleton. Thus it is closed, and its subtype inclusion is a
closed embedding. No separation property of the total scheme is needed.
\end{proof}

\begin{theorem}[Continuous native morphisms over the actual base]
\label{scheme:relative-valued-point-functor}
\lean{EnvelopeProperty.Berkovich.RelativeSchemeValuedPoint.map,
EnvelopeProperty.Berkovich.RelativeSchemeValuedPoint.map_coe,
EnvelopeProperty.Berkovich.RelativeSchemeValuedPoint.continuous_map,
EnvelopeProperty.Berkovich.RelativeSchemeValuedPoint.map_id,
EnvelopeProperty.Berkovich.RelativeSchemeValuedPoint.map_comp,
EnvelopeProperty.Berkovich.relativeSchemeValuedPointFunctor}
\leanok
\uses{scheme:relative-valued-points,scheme:valued-point-pullback}
A commuting scheme morphism over $\operatorname{Spec}R$ induces a continuous
map of the corresponding native prescribed-base spaces. The maps obey whole
point identity and composition laws and form an actual functor from schemes
over $\operatorname{Spec}R$ to topological spaces.
\end{theorem}
\begin{proof}
Native composition identifies the image under the target structure morphism
with the image under the source structure morphism. The commuting diagram
therefore preserves the exact fixed base point. Restrict the continuous native
map to the two subspaces. Whole-point identity and composition descend from
the genuine residue-field pullback laws. These laws give the actual over-category
functor, without any point-existence or scalar-extension-surjectivity assumption.
\end{proof}

\begin{theorem}[Actual affine prescribed-base comparison]
\label{scheme:relative-affine-valued-point-homeomorphism}
\lean{EnvelopeProperty.Berkovich.RelativeSchemeValuedPoint.mem_affine_relativeLocus_iff,
EnvelopeProperty.Berkovich.RelativeSchemeValuedPoint.affineSpectrumHomeomorph,
EnvelopeProperty.Berkovich.RelativeSchemeValuedPoint.affineSpectrumHomeomorph_apply_coe,
EnvelopeProperty.Berkovich.RelativeSchemeValuedPoint.affineSpectrumHomeomorph_symm_apply_coe,
EnvelopeProperty.Berkovich.RelativeSchemeValuedPoint.affineStructureMap_commutes,
EnvelopeProperty.Berkovich.RelativeSchemeValuedPoint.affineSpectrumHomeomorph_map}
\leanok
\uses{scheme:relative-valued-points,scheme:relative-valued-point-functor,
scheme:affine-valued-point-homeomorphism,spectrum:relative}
For an actual $R$-algebra $A$, the native coefficient locus for the actual
map $\operatorname{Spec}A\to\operatorname{Spec}R$ is homeomorphic to
$\mathcal S_v(A)$. It preserves every element value and commutes with
the genuine Spec maps of $R$-algebra homomorphisms.
\end{theorem}
\begin{proof}
The native affine homeomorphism is natural for the actual coefficient ring
homomorphism. Mapping a native point to the fixed base point is therefore
equivalent to restricting its entire original-ring seminorm to $v$.
Equality of the whole seminorm means equality on every coefficient, precisely
the defining relative-spectrum locus. Restrict the genuine absolute
homeomorphism to this exact pair of subspaces. The whole-point ring-map
naturality gives element preservation and relative algebra-map naturality.
No field, triviality, completeness or finite-type hypothesis is used.
\end{proof}

\begin{theorem}[Genuine coefficient action on scheme section rings]
\label{scheme:relative-section-coefficients}
\lean{EnvelopeProperty.Berkovich.SchemeValuedPoint.sectionBaseMap,
EnvelopeProperty.Berkovich.SchemeValuedPoint.sectionBaseAlgebra,
EnvelopeProperty.Berkovich.SchemeValuedPoint.sectionBaseMap_apply,
EnvelopeProperty.Berkovich.RelativeSchemeValuedPoint.eval_sectionBaseMap,
EnvelopeProperty.Berkovich.RelativeSchemeValuedPoint.sectionSpectrum,
EnvelopeProperty.Berkovich.RelativeSchemeValuedPoint.sectionSpectrum_val_apply,
EnvelopeProperty.Berkovich.RelativeSchemeValuedPoint.continuous_sectionSpectrum}
\leanok
\uses{scheme:relative-valued-points,scheme:valued-point-pullback,scheme:valued-points,
scheme:valued-point-topology,spectrum:relative}
For a genuine structure morphism $f:X\to\operatorname{Spec}R$ and any
scheme open $U$, compose the actual global section identification of $R$,
the actual structure-map section pullback, and genuine restriction to $U$.
This defines a canonical coefficient ring map $R\to\Gamma(X,U)$ and its
induced $R$-algebra action. Evaluating these actual coefficients at every
native prescribed-base point supported in $U$ gives exactly $v$. The
resulting actual relative-spectrum section evaluation is continuous on
its genuine local domain.
\end{theorem}
\begin{proof}
The composite is a ring homomorphism because all three maps are genuine
scheme section ring maps. It supplies the algebra action rather than assuming
an unrelated one. Actual section restriction preserves residue evaluation,
and actual global evaluation naturality identifies its coefficient value
with evaluation at the mapped native base point. That whole point is the
fixed inverse image of $v$, so the affine evaluation identity gives $v(r)$
for every $r\in R$. Thus genuine local section evaluations also define
relative-spectrum points for this actual coefficient action. Their value
on every local section is the original native evaluation. Every such
evaluation is continuous on the open domain where the section is defined;
compose with the two subtype inclusions and use the coordinate criterion
for relative-spectrum continuity.
\end{proof}
